\documentclass[11pt]{letter}
\usepackage[T1]{fontenc}
\usepackage{newtxtext}
\usepackage{newtxmath}
\usepackage[margin=1in]{geometry}
\usepackage{xurl}

\signature{Yuan Qiu\\
Department of Statistics, College of Economics\\
Hangzhou Dianzi University\\
Hangzhou 310003, China\\
\texttt{yuanqiu@hdu.edu.cn}}
\address{Yuan Qiu\\
Department of Statistics, College of Economics\\
Hangzhou Dianzi University\\
Hangzhou 310003, China\\
\texttt{yuanqiu@hdu.edu.cn}}

\begin{document}

\begin{letter}{Editor\\ \emph{Ecological Informatics}}

\opening{Dear Editor,}

I am pleased to submit my manuscript entitled ``Validation and Calibration Diagnostics for Migratory-Route Comparison from Incomplete Animal Tracking Data'' for consideration in \emph{Ecological Informatics}.

This manuscript presents a reproducible diagnostic framework for validating, calibrating, and propagating missing-route uncertainty in migratory-route comparison. Using open GPS tracking data for Lesser Black-backed Gulls from the \texttt{LBBG\_ZEEBRUGGE} dataset, the study compares random point loss with contiguous tracking gaps, evaluates deterministic spherical bridge reconstruction, validates Brownian bridge uncertainty envelopes against withheld route segments, estimates empirical radius inflation needed for nominal coverage, and propagates reconstruction uncertainty into distance matrices, clustering, anomaly ranking, and top-ranked anomaly overlap.

The paper is positioned as an ecological-informatics methods contribution. Its main claim is that route-comparison conclusions should not rely only on a single deterministic distance matrix when prolonged tracking gaps are present. Instead, plausible reconstruction uncertainty should be validated and calibrated where possible, then propagated into the downstream conclusions that ecologists inspect.

The empirical results are intentionally framed as a reproducible open-data case study rather than as species-general biological claims about gull migration. The transferable contribution is the diagnostic framework for validating, calibrating, and propagating missing-route uncertainty into downstream computational conclusions.

I believe the manuscript fits \emph{Ecological Informatics} because it addresses a practical computational problem in ecological tracking data analysis and provides a transparent, reproducible framework for uncertainty-aware movement data comparison.

This manuscript is original, has not been published previously, and is not under consideration elsewhere.

The data used in the study are publicly available from the INBO \texttt{LBBG\_ZEEBRUGGE} dataset. The analysis code and manuscript materials are available at:

\url{https://github.com/cliffyb824/bird-migration-trajectory-metrics}

The author declares no competing interests. This research did not receive any specific grant from funding agencies in the public, commercial, or not-for-profit sectors.

\closing{Sincerely,}

\end{letter}

\end{document}
